% Claim statement file for row claim_3_20 -- "AcyclicNonCollidersBlockable".
% Auto-generated stub by scaffold/solve_chapter.py at row start. The
% manager agent (or a worker it dispatches) fills the body below with the
% claim's statement(s) from the lecture notes -- statement only, no proof
% (proofs live in the sibling `claim_3_20_proof_*.tex` file).
%
% This file is a sibling of `main.tex` inside `<subsection>/tex/`. The
% `subfiles` package lets it render both standalone and as part of
% `main.tex`. Note: `main.tex` does NOT \subfile this statement file -- the
% sibling `_proof_` file restates the statement above the proof, so
% including both in `main.tex` would be redundant. This file exists for
% standalone reading / grep convenience.

\documentclass[main]{subfiles}

\begin{document}
% ref: claim_3_20 (statement)
% title: AcyclicNonCollidersBlockable
\def\rowref{claim\_3\_20}
\phantomsection\label{claim_3_20}

\begin{Rem}[AcyclicNonCollidersBlockable]
    Let $G = (J, V, E, L)$ be a CDMG (def \ref{def-cdmg}) and suppose $G$ is acyclic in the sense of def \ref{def-acylic}. Then, for every integer $n \ge 0$, every walk
    \[
        \pi = (v_0, a_0, v_1, a_1, \dots, v_{n-1}, a_{n-1}, v_n)
    \]
    in $G$ in the sense of def \ref{def:walks}, item~i.\ (so $v_0, \dots, v_n \in J \cup V$, $a_0, \dots, a_{n-1} \in E \cup L$, and the walk constraint
    \[
        \bigl(a_k = (v_k, v_{k+1}) \;\land\; a_k \in E \cup L\bigr)
        \;\lor\;
        \bigl(a_k = (v_{k+1}, v_k) \;\land\; a_k \in E\bigr)
    \]
    holds at every $k \in \{0, 1, \dots, n - 1\}$), and every position $k \in \{0, 1, \dots, n\}$, the following implication holds:
    \[
        \bigl(\text{$k$ is a non-collider on $\pi$ per def \ref{def:collider_noncollider}, item~i.}\bigr)
        \;\implies\;
        \bigl(\text{$k$ is a blockable non-collider on $\pi$ per def \ref{def:unblockable_noncollider}}\bigr).
    \]
\end{Rem}

\end{document}
